\emph{Conventions.} Fermions are two-component, $\kappa = \tfrac12$, and Yukawa couplings are written $-\mathcal{L} \supset \tfrac12 Y^{a}_{jk}\psi_{j}\psi_{k}\phi_{a} + \mathrm{h.c.}$ with real scalars $\phi_{a}$, as in Fact \claimref{synthesis.imported-two-loop-yukawa-gauge-terms}. A complex scalar with components $\Phi_{c} = (\phi_{c1} + i\phi_{c2})/\sqrt{2}$ along an orthonormal basis $T^{c}$ of its place in the fermion bilinear, coupled as $\tfrac12(y\otimes T^{c})_{jk}\psi_{j}\psi_{k}\Phi_{c}$, has $Y^{c1} = (y\otimes T^{c})/\sqrt2$ and $Y^{c2} = i(y\otimes T^{c})/\sqrt2$, so $\sum_{a}Y^{a}Y^{a\dagger} = yy^{\dagger}\otimes\sum_{c}T^{c}T^{c\dagger}$. For the $(10,1,3)$ the entry along $\nu^{c}\nu^{c}$ is a basis element, so $\langle\Delta_{R}\rangle = v_{R}$ gives $\tfrac12 v_{R}f\,\nu^{c}\nu^{c}$, that is $M_{R} = v_{R}f$ as in Fact \claimref{standard-model.imported-clebsch-structure-seesaw}. A Dirac-type coupling $(y\otimes T^{c})_{(iA)(jB)}\psi_{iA}\psi^{c}_{jB}\Phi_{c}$ fills the two off-diagonal blocks of $Y^{a}$, and $\sum_{a}Y^{a}Y^{a\dagger}$ is $yy^{\dagger}\otimes\sum_{c}T^{c}T^{c\dagger}$ on the $\psi$ block and $y^{\mathsf T}\bar{y}\otimes\sum_{c}T^{c\mathsf{T}}\bar{T}^{c}$ on the $\psi^{c}$ block.

\emph{(i)} Each of $\psi = (4,2,1)$ and $\psi^{c} = (\bar4,1,2)$ has eight components $(a, \alpha)$. Since $\Phi$ spans an irreducible representation and $\psi$, $\psi^{c}$ are irreducible, $\sum_{c}T^{c}T^{c\dagger}$ is invariant on each eight-dimensional space and so, by Schur's lemma, $(\dim R/8)\,\mathbb{1}_{8}$, its trace being $\sum_{c}\mathrm{Tr}\,T^{c}T^{c\dagger} = \dim R$. The check confirms it on explicit orthogonal bases: $30/8$ for the $(10,1,3)$ (symmetric in the $SU(4)$ pair and in the $SU(2)_{R}$ pair), $4/8$ for the $(1,2,2)$ ($\delta_{ab}\otimes E_{\alpha\beta}$), $60/8$ for the $(15,2,2)$ (off-diagonal $E_{ab}$ and traceless diagonals), on both blocks. With $C_{2}(4) = \tfrac{15}{8}$, $C_{2}(2) = \tfrac34$, $d(SU(4)) = 15$, $d(SU(2)) = 3$:
\[
Y_{4}^{i} = \frac{1}{d(G_{i})}\mathrm{Tr}\bigl[C_{2}^{i}(F)\,Y^{a}Y^{a\dagger}\bigr] = \frac{\dim R\,\mathrm{Tr}(yy^{\dagger})}{d(G_{i})}\bigl(C_{2}^{i}(\psi) + C_{2}^{i}(\psi^{c})\bigr),
\]
which is $\dim R\,\mathrm{Tr}(yy^{\dagger})\cdot(\tfrac{15/4}{15}, \tfrac{3/4}{3}, \tfrac{3/4}{3}) = \tfrac14\dim R\,\mathrm{Tr}(yy^{\dagger})(1,1,1)$, because $\psi$ and $\psi^{c}$ both carry the fundamental of $SU(4)$ and each carries the doublet of exactly one $SU(2)$. The matching gives $(\tfrac35 + \tfrac25)\cdot\tfrac14\dim R\,\mathrm{Tr}(yy^{\dagger}) = \tfrac14\dim R\,\mathrm{Tr}(yy^{\dagger})$ for $a_{1}$ and the same for $a_{2}$, $a_{3}$: a universal vector. For the $(10,1,3)$ only $\psi^{c}$ couples, $\mathrm{Tr}\,YY^{\dagger}$ over $\psi^{c}$ is $8\cdot\tfrac{30}{8}\mathrm{Tr}(ff^{\dagger}) = 30\,\mathrm{Tr}(ff^{\dagger})$, so $Y_{4} = (30\cdot\tfrac{15}{8}/15,\ 0,\ 30\cdot\tfrac34/3)\mathrm{Tr}(ff^{\dagger}) = (\tfrac{15}{4}, 0, \tfrac{15}{2})\mathrm{Tr}(ff^{\dagger})$, matched $(\tfrac35\cdot\tfrac{15}{2} + \tfrac25\cdot\tfrac{15}{4},\ 0,\ \tfrac{15}{4}) = (6, 0, \tfrac{15}{4})$. The $(10,1,3)$ cannot couple as $\psi\psi^{c}$ ($4\otimes\bar4$ holds no $10$), and by Fact \claimref{standard-model.imported-pati-salam-representation-theory-survival}(b) and Proposition \claimref{standard-model.majorana-element} the only scalars light on $(M_{I}, M_{U})$ in the chains are the bidoublet, the $(10,1,3)$ and, where named, the $(15,2,2)$; the $(\overline{10},3,1)$ is not light.

\emph{(ii)} By the Fact, the Yukawa part of the gauge running is $d\alpha_{i}^{-1}/d\ln\mu \supset \kappa Y_{4}^{i}/16\pi^{3}$ (from $\beta(g) \supset -2\kappa g^{3}Y_{4}/(4\pi)^{4}$ and $d\alpha^{-1}/d\ln\mu = -8\pi\beta/g^{3}$); integrated over $\ln(M_{U}/M_{I}) = 2\pi y$ it is $\kappa\langle Y_{4}^{i}\rangle y/8\pi^{2}$, which in $a_{i}(M_{Z}) = a_{U} + b^{\mathrm{SM}}_{i}x + b^{\mathrm{eff}}_{i}y$ enters, after the matching of Lemma \claimref{standard-model.matching-traces} at $M_{I}$, as $-\varepsilon v_{i}y - u$ with $\varepsilon = \langle\mathrm{Tr}(ff^{\dagger})\rangle/16\pi^{2}$, $v = (6, 0, \tfrac{15}{4})$ and $u$ universal. (The sign of $u$ is immaterial.) Here $b^{\mathrm{SM}}$ and $b^{\mathrm{eff}}$ are the one-loop sums of Fact \claimref{standard-model.imported-coefficient-sum-indices} and the $a_{i}(M_{Z})$ the inputs of Proposition \claimref{standard-model.spectrum-run}. For fixed $\varepsilon$ and $u$ the system is linear in $(a_{U}, x, y)$ with matrix rows $(1, b^{\mathrm{SM}}_{i}, b^{\mathrm{eff}}_{i} - \varepsilon v_{i})$; its determinant is affine in $\varepsilon$, computed as $D_{0} + \tfrac{201}{4}\varepsilon$ with $D_{0} = \tfrac{364}{15}, \tfrac{386}{15}, \tfrac{114}{5}$ for the three chains, positive for every $\varepsilon \ge 0$; expanding along the third column, the coefficient of $\varepsilon$ is $-\det[\mathbb{1}, b^{\mathrm{SM}}, v] = -[6(b^{\mathrm{SM}}_{3} - b^{\mathrm{SM}}_{2}) + \tfrac{15}{4}(b^{\mathrm{SM}}_{2} - b^{\mathrm{SM}}_{1})] = \tfrac{201}{4}$, which involves only the columns $\mathbb{1}$ and $b^{\mathrm{SM}}$ and so is the same for every chain (checked directly). So Cramer's rule gives one solution for every $\varepsilon \ge 0$; it back-substitutes, and the check verifies on $\varepsilon \in \{0, \tfrac18, \dots, 3\}$ (the cap $\varepsilon \le 3$ is singular values of $f$ at most $4\pi$) that $x$ rises strictly, and that a universal $u$ shifts $a_{U}$ by $-u$ and leaves $x$, $y$ unchanged. Three data, the measured $a_{i}(M_{Z})$, are used up by the three unknowns $(a_{U}, x, y)$ at each $\varepsilon$; no equation remains for $\varepsilon$. The two-loop gauge terms of Proposition \claimref{standard-model.chain}(1) add to this system terms of relative order $\alpha/4\pi$ that are smooth in $(a_{U}, x, y)$ and independent of $\varepsilon$; the Jacobian stays nonsingular, so by the implicit function theorem the same conclusion holds for the full two-loop system near the stored solution (argued, not checked). The running of $f$ itself only enters through the average $\langle\mathrm{Tr}(ff^{\dagger})\rangle$, so no further datum appears.

\emph{(iii)} $x(\varepsilon)$ is a ratio of affine functions; the check extracts $dx/d\varepsilon$ at $\varepsilon = 0$ exactly ($\tfrac{2373195393}{977158000}$, $\tfrac{605922228}{274711375}$, $\tfrac{129038993}{47922750}$) and, with $d\ln M_{I} = 2\pi\,dx$ and $\varepsilon = \mathrm{Tr}/16\pi^{2}$, $d\ln M_{I}/d\,\mathrm{Tr}(ff^{\dagger}) = (dx/d\varepsilon)/8\pi$, bounded with certified rational bounds on $\pi$ into $(0.09, 0.10)$, $(0.08, 0.09)$, $(0.10, 0.11)$. Moving $\alpha_{s}$ by $\pm 0.0009$ from the input $59/500$ moves $\ln M_{I}$ by $\mp(0.03$ to $0.05)$ (computed: $-0.0365/+0.0371$, $-0.0334/+0.0339$, $-0.0401/+0.0407$). If $x$ were supplied, the system becomes linear in $(a_{U}, y, z = \varepsilon y)$; solving it at the $x$ of $\varepsilon = 1/158$ returns $\varepsilon = 1/158$ exactly, and with $\alpha_{s}$ moved by $\pm 0.0009$ returns traces differing from it by more than $0.3$ (computed: $1.37$ and $0.61$ for the real chain, similar for the others). For comparison (computed, not checked): moving the $(10,1,3)$ threshold from $M_{I}$ to $1.1\,M_{I}$ moves $\ln M_{I}$ by $0.017$ to $0.027$, so a unit trace matches a threshold displaced by a factor $1.5$ to $1.8$, well within the unknown spread of the heavy masses.

\emph{The orbit.} Under $g_{\lambda}$ of Theorem \claimref{synthesis.majorana-scale-free-on-pati-salam-route}, $\varepsilon \mapsto |\lambda|^{2}\varepsilon$, so $M_{I}$ and $v_{R} = M_{I}/g_{R}$ rise strictly with $|\lambda|$ by (ii), while (EW) and the Dirac matrices stay fixed; $m_{3} = m_{D}^{2}/(f_{33}v_{R})$ scales as $1/(|\lambda|\,v_{R}(|\lambda|))$, strictly decreasing, so the observed $m_{3}$ fixes $|\lambda|$ uniquely and the Majorana scale remains valued only by $m_{3}$ (the slow variation of $g_{R}(M_{I})$ is of the same order as the shift and does not reverse the monotonicity at these sizes; argued). The numbers in (iii) use the one-loop solutions as baseline; the stored two-loop scales differ, which changes the sizes at the tens-of-percent level and not the count.
